%
%****************************************************************************************************
%======================= B E G I N   D O C U M E N T   S E T T I N G S ========================
%****************************************************************************************************
%
\documentclass[a4paper, 12pt, xcolor=dvipsnames]{scrartcl}	% Article als Modul
\usepackage[ngerman]{babel}
\usepackage[onehalfspacing]{setspace}
\usepackage{fancyhdr}
\usepackage[usenames, ps2pdf, dvipsnames,svgnames,table]{xcolor}
\usepackage[utf8x]{inputenc}
\usepackage{wrapfig}   
\usepackage{color}
\usepackage[printonlyused, withpage]{acronym}
\usepackage{ifthen}
\usepackage{lastpage}
\usepackage{currfile}
\usepackage{caption}						% Formatierung der Bildunterschriften
\usepackage{sidecap}
\usepackage{subcaption}
\usepackage{tikz}
\usepackage{calc}
\usepackage{graphicx}
\usepackage{multirow}
\usepackage{longtable}
\usepackage{amsthm}
\usepackage{amsbsy}
\usepackage{amssymb}
\usepackage{array}
\usepackage{chngpage}
\usepackage{tabulary} 					% instead of tabularx
\usepackage{array,longtable,calc,amsmath}
\usepackage{xspace}
\usepackage{pdfpages}
\usepackage{textfit}
\usepackage{tabu}
\usepackage{setspace}
\usepackage{listings}
\pagestyle{fancy}
\usepackage[acronym,toc]{glossaries}
\usepackage{docmute}					% includes parts of files just within the Begin and End Document
\usepackage{blindtext}
\usepackage{lipsum}
%---------------------Switch modulare Printversionen
%
\newboolean{UseWm} 		%Korrektur Wasserzeichen in das PDF drucken
\setboolean{UseWm}{false} 		%Zuweisung eines Wertes                    			% Input file für gemeinsame Package Definitionen
%
%--------------------Set data for document:
\newcommand{\SYear}{2021/22}
\newcommand{\SClass}{5AHET}
\newcommand{\ADatum}{04. April 2022}
\newcommand{\DAName}{E$^2$MILY - Elektrische Hardwaredokumentation}
\newcommand{\DPLNameOnePre}{Vorname1}						% Diplomand 1 als Variablenname, usage: \DPLName...
\newcommand{\DPLNameOneSur}{Nachname1}					% Diplomand 1 als Variablenname, usage: \DPLName...
\newcommand{\DPLNameTwoPre}{Vorname2}						% Diplomand 2 als Variablenname, usage: \DPLName...
\newcommand{\DPLNameTwoSur}{Nachname2}					% Diplomand 2 als Variablenname, usage: \DPLName...
\newcommand{\DPLNameThreePre}{Vorname3}					% Diplomand 3 als Variablenname, usage: \DPLName...
\newcommand{\DPLNameThreeSur}{Nachname3}				% Diplomand 3 als Variablenname, usage: \DPLName...
\newcommand{\DPLNameFourPre}{Vorname4}					% Diplomand 4 als Variablenname, usage: \DPLName...
\newcommand{\DPLNameFourSur}{Nachname4}					% Diplomand 4 als Variablenname, usage: \DPLName...
\newcommand{\BNameOnePre}{Vorname1}							% Betreuer 1 als Variablenname, usage: \DPLName...
\newcommand{\BNameOneSur}{Nachname1}						% Betreuer 1 als Variablenname, usage: \DPLName...
\newcommand{\BNameTwoPre}{Vorname2}						% Betreuer 2 als Variablenname, usage: \DPLName...
\newcommand{\BNameTwoSur}{Nachname2}						% Betreuer 2 als Variablenname, usage: \DPLName...
%
\setboolean{UseWm}{true} 												% Wasserzeichen einschalten
%\newcommand{\WatermarkName}{Korrekturexemplar}		%Text für das Wasserzeichen
%\newcommand{\WatermarkName}{01.01.2021}					%AlternativText für das Wasserzeichen
\newcommand{\WatermarkName}{Template}						%AlternativText für das Wasserzeichen
%
%==============================S E T T I N G S ===================================
%
\definecolor{green}{rgb}{0.8,1.0,0.2}
\definecolor{orange}{rgb}{0.9,0.64,0.25}
\definecolor{kaki}{rgb}{0.74,0.74,0.088}
\definecolor{blue}{rgb}{0.2,0.59,0.94}
\definecolor{dred}{rgb}{0.722,0.18,0.18}
\definecolor{htl}{rgb}{0.15625,0.4375,0.507}
%
%------------------------------Watermark wird mit einer Designvariablen geladen ---------------------------------------------------------------------------------
%
\ifthenelse{\boolean{UseWm}}%
{\usepackage{draftwatermark}\SetWatermarkText{\WatermarkName$~-~$\WatermarkName\xspace$~-~$\WatermarkName} \SetWatermarkLightness{0.95}\SetWatermarkScale{0.5}}{}
%
\graphicspath{{./../03_AUX/01_PICT/}}
\newcommand{\changefont}[3]{\fontfamily{#1} \fontseries{#2} \fontshape{#3} \selectfont} % Change a new font
%
% --------------------Set the figure description:
%
\captionsetup{margin=10pt,font=small, labelfont={color=black!80, bf}, textfont={color=black!60}, format=hang, indention=-1cm}
\captionsetup[wrapfigure]{name=Bild}
\captionsetup[figure]{name=Bild}
% --------------------Set the page margins:
%
\setlength{\voffset}{1.0cm}
\setlength{\textwidth}{17.0cm}
%\addtolength{\textheight}{10.0cm}
\setlength{\textheight}{24.0cm}
\setlength{\topmargin }{-2.3cm}
\setlength{\marginparsep }{0pt}
\setlength{\headsep }{1.2cm}
\setlength{\oddsidemargin}{0.0cm}
\setlength{\evensidemargin}{0.0cm}
\setlength{\marginparwidth}{0pt}
\setlength{\hoffset}{-0.54cm}
\setlength{\headwidth}{17.0cm}
\setlength{\footskip}{30pt}
\setlength{\parindent}{0pt}
%
%

%--------------------Set environment for Header and and Footer:
%
\pagestyle{fancy}
\rhead{\includegraphics[scale=0.06]{/HTL_logo.pdf}}
\chead{}
\lhead{\begin{tabular}{p{12.0cm}} 
\small{\changefont{pag}{b}{sc}\textcolor{black!30}{\DAName}}
\end{tabular} \vspace{0.1cm}}
\lfoot{\small{\changefont{pag}{b}{sc}\textcolor{black!30}{\SYear - \SClass - \DPLNameOneSur, \DPLNameTwoSur, \DPLNameThreeSur, \DPLNameFourSur}}} 
\cfoot{}
\rfoot{\small{\changefont{pag}{b}{sc}\textcolor{black!30}{\thepage\ \textbf{\textcolor{htl!100}I} \pageref{LastPage}}}}
\renewcommand{\headrulewidth}{0.4pt}
\renewcommand{\footrulewidth}{0.4pt}
\let\myHeadrule\headrule
\let\myFootrule\footrule
\renewcommand\headrule{\color{black!30}\myHeadrule }
\renewcommand\footrule{\color{black!30}\myFootrule }
\addto{\captionsngerman}{%
  \renewcommand*{\contentsname}{Inhalt}
  \renewcommand*{\listfigurename}{Abbildungen}
  \renewcommand*{\listtablename}{Abbildungsverzeichnis}
}
%\pagestyle{fancy}
\fancypagestyle{ErsteSeite}{%
   \fancyhf{}%
   \fancyhead[L]{}
}

%\makenoidxglossaries
\renewcommand*\acronymname{Vokabelverzeichnis}
\renewcommand*\glossaryname{Vokabelverzeichnis}

%--------------------Neue Umgebung für M E R K S A T Z:
%
%-------------------- Zählerdefinition:
%
\newcounter{MsNr}
\renewcommand{\theMsNr}{\arabic{MsNr}}
%
\newlength{\boxw}
\newlength{\boxh}
\newlength{\shadowsize}
\newlength{\boxroundness}
\newlength{\tmpa}
\newsavebox{\shadowblockbox}
\setlength{\shadowsize}{6pt}
\setlength{\boxroundness}{3pt}
%
\newenvironment{shadowblock}[1]%
{\begin{lrbox}{\shadowblockbox}\begin{minipage}{#1}}%
{\end{minipage}\end{lrbox}%
\settowidth{\boxw}{\usebox{\shadowblockbox}}%
\settodepth{\tmpa}{\usebox{\shadowblockbox}}%
\settoheight{\boxh}{\usebox{\shadowblockbox}}%
\addtolength{\boxh}{\tmpa}%
\begin{tikzpicture}
\addtolength{\boxw}{\boxroundness * 2}
\addtolength{\boxh}{\boxroundness * 2}
\foreach \x in {0,.05,...,1}
{
\setlength{\tmpa}{\shadowsize * \real{\x}}
\fill[xshift=\shadowsize - 1pt,yshift=-\shadowsize + 
1pt,black,opacity=.04,rounded corners=\boxroundness] 
(\tmpa, \tmpa) rectangle +(\boxw+20 - \tmpa - \tmpa, \boxh+20 - \tmpa - 
\tmpa);
}
\filldraw[fill=yellow!10, draw=kaki!100, rounded corners=\boxroundness] (0, 
0) rectangle (\boxw+20, \boxh+20);
\draw node[xshift=\boxroundness,yshift=\boxroundness,inner sep=10pt,outer 
sep=0pt,anchor=south west] (0,0) {\usebox{\shadowblockbox}};
\end{tikzpicture}}
%
%====================================================================
\newenvironment{Merksatz}
{\begin{shadowblock}{15.8cm}
{\par\vspace*{0.0cm}
\noindent \refstepcounter{MsNr}
\textcolor{kaki!100} {\large{}\textbf{Merke \theMsNr :}} % Hier kann man statt Merke auch etwas anderes hinschreiben
\label{Msl-a}}
}
{\end{shadowblock}}
%
%--------------------End    M E R K S A T Z
% Settings for Codelistings within documentation:
\lstset{ % General setup for the package
	language=C,
	basicstyle=\small\sffamily,
	numbers=left,
 	numberstyle=\tiny,
	frame=tb,
	tabsize=4,
	columns=fixed,
	showstringspaces=false,
	showtabs=false,
	keepspaces,
	commentstyle=\color{red},
	keywordstyle=\color{blue}
}

\setcounter{footnote}{0}

\changefont{pag}{m}{n}                    				% Input file für gemeinsame dokumentenbezogene Definitionen
%
\begin{document}																%ab hier wird der Text im Hauptdokument übernommen
% Verzeichnisse (to be commented out for compilation of main file):
%
%\tableofcontents 																%ACHTUNG: für die Gesamtcompilierung auskommentieren!!
%\newpage 																		%ACHTUNG: für die Gesamtcompilierung auskommentieren!!
%
\changefont{pag}{m}{n}
%
%****************************************************************************************************
%======================== B E G I N  S U B  D O C U M E N T  ==========================
%****************************************************************************************************
%
\section{Projektmanagement}
Um den zeitlichen Ablauf der Diplomarbeit zu planen und zu dokumentieren sind einigen Dokumente und Programme erforderlich.
\begin{enumerate}																%Beginn Nummerierung
	\item\textbf{WORD File:} Erstellen Sie ein Pflichtenheft laut Vorlage!
	\item\textbf{EXCEL File:} Führen Sie eine Art Tagebuch über den Verlauf der DA mit den eingetragenen Meilensteinen.
\end{enumerate}

\section{Dokumentation der Diplomarbeit}							%Hauptüberschrift
Erstellen Sie die Dokumentation entsprechend der Vorlagen der HTL-Weiz.\\
\textbf{Wieder gilt:} Ein früher Beginn sichert den Erfolg! Folgende Punkte sind wichtig:
\begin{enumerate}																%Beginn Nummerierung
	\item\textbf{Dokumentstruktur:} Doppelseitig					% Aufzählung Eintrag
	\item\textbf{Schriftart:} Arial 12										% Aufzählung Eintrag
	\item\textbf{Zeilenabstand:} 1.15										% Aufzählung Eintrag
\end{enumerate}

\vspace{-0.6cm}																	% vertikaler Abstand

\subsection{Zeichnungen}													%Subüberschrift
Ich empfehle die Zeichnungen, die für die Dokumentation der DA notwendig sind, selbst und mit der Software \textbf{VISIO} zu zeichnen. Das Programm ist über das Office Paket für Sie gratis.\\
Beginnen Sie schon früh mit dem Erstellen solcher Blockschaltbilder und technischer Zeichnungen, dann geht das Dokumentieren am Ende der DA umso schneller.\\
\vspace{-0.6cm}																	% vertikaler Abstand

\pagebreak

\section{Konstruktion und Entwurf von Baugruppen}				%Hauptüberschrift
\subsection{Berechnungen}													%Subüberschrift
Führen Sie die für die \acs{DA} (hier steht in der Source nur die Abkürzung DA) notwendigen Berechnungen mit dem Programm \textbf{sMath} oder \textbf{SCILAB} durch. Beginnen Sie hier ebenfalls so früh wie möglich, um darin Routine zu bekommen. Eine Beschreibung für die Einführung in sMath und SCILAB ist vorhanden.
\vspace{-0.6cm}																	% vertikaler Abstand

\subsection{Simulation elektronischer Schaltungen}				%Subüberschrift
Ein kostenloses und sehr flexibles Programm Simulationen elektrischer Schaltungen durchführen zu können ist \textbf{LT-Spice} von Linear Technology. Eine Beschreibung für die Einführung in die Simulationstechnik ist im Internet vorhanden.
\vspace{-0.6cm}																	% vertikaler Abstand

\subsection{Platinenlayout}													%Subüberschrift
Für die professionelle und automatische Erstellung von Layouts wird das Programm \textbf{KiCad} verwendet. Eine Dokumentation ist ebenfalls vorhanden.\\
Beachten Sie beim Entwurf des Layouts einer Schaltung Dokumente und Richtlinien, die Sie von mir bekommen. Darin werden Hinweise zum Entwurf, sowie das \textbf{EMV} und \textbf{ESD}-gerechte Leiterplattendesign behandelt.

\section{Kommunikation}													%Hauptüberschrift
\subsection{Schriftliche Korrespondenz}								%Subüberschrift
Die schriftliche Kommunikation wird hauptsächlich über email oder TEAMS geführt. Beachten Sie folgende Punkte:
\begin{enumerate}																%Beginn Nummerierung
\item\textbf{Zieladresse:} Bitte senden Sie Ihre mails immer an: \textbf{eBergler@lr.htlweiz.at} - oftmals schreibe ich von anderen Adressen zurück. Verwenden Sie jedoch immer die HTL Addresse! Nur so ist es gewährleistet, dass ich die Nachricht so schnell als möglich bekomme.
\item\textbf{Addressaten:} Schreiben Sie Ihre Nachricht immer an \textbf{ALLE} die an Ihrer DA beteiligt sind, damit stellen Sie sicher, dass jeder den gleichen Informationsstand hat. Meine Telefonnummer lautet: \textbf{0664 882 5 9797}.
\end{enumerate}
\textbf{Grundsätzlich:} Es werden keine Daten, Dokumente oder Files über email versendet $\to$ dazu gibt es Laufwerke in einer Cloud!\\
 

\subsection{Datenaustausch}													%Subüberschrift
Elektronische Daten werden \textbf{ausschließlich} über eine Ordnerstruktur in einer Cloud ausgetauscht. Dies dient der Sicherung der Datenkonsistenz und Übersichtlichkeit. 

\section{Beurteilung der Diplomarbeit}										%Hauptüberschrift
\subsection{Schriftliche und praktische Beurteilungskriterien der Arbeit}	%Subüberschrift
\begin{enumerate}																	%Beginn Nummerierung
	\item\textbf{Fachkompetenz:} Wie gut kennt sich der Diplomand im Fachgebiet aus!
	\item\textbf{Methodenkompetenz:} Welche Methoden zur Problemlösung wurden gewählt!
	\item\textbf{Selbstkompetenz:} Hier fließt Ihre Motivation und Termintreue ein!
	\item\textbf{Sprachkompetenz:} Wie gut kann der Diplomand technische Zusammenhänge sprachlich richtig darstellen!
	\item\textbf{Dokumentation:} Wie gut ist die Dokumentation der Arbeit ausgeführt!
\end{enumerate}

\subsection{Mündliche Beurteilungskriterien der Präsentation und Verteidigung}		%Subüberschrift
\begin{enumerate}																	%Beginn Nummerierung
	\item\textbf{Zeitliche Einteilung:} Sie können Ihre Diplomarbeit im vorgegebenen Zeitrahmen darstellen.
	\item\textbf{Inhalt:} Die inhaltlich-fachliche Auseinandersetzung mit dem Thema ist in der erforderlichen Tiefe gegeben.
	\item\textbf{Argumentation:} Sie können Ihre Standpunkte überzeugend argumentieren.
	\item\textbf{Rethorik:} Der Kandidat kann seine Diplomarbeit in Standardsprache klar strukturiert und verständlich darstellen.
	\item\textbf{Medienumgang:} Sie können mit zeitgemäßen Präsentationsmitteln umgehen.
\end{enumerate}

%----------------------------------------
\Blindtext \lipsum[1-2]															%Blindtext zum Auffüllen mit unsinnigen Sätzen
%----------------------------------------
%
\newpage																			%Seitenumbruch am Ende des Subkapitels
%****************************************************************************************************
%============================= E N D   S U B D O C U M E N T  ==========================
%****************************************************************************************************
%
% Verzeichnisse (to be commented out for compilation of main file):
%
%% Zusammenfassung und Resumee der Diplomarbeit:
\section{Zusammenfassung}

Hier steht die Zusammenfassung und das Resumee der Diplomarbeit!
\newpage

% Anhang zur Diplomarbeit:
\section{Anhang}

Hier stehen zusätzliche Informationen zur Diplomarbeit!
Zum Beispiel können hier Listings oder Teile des Programmes abgebildet sein:

\begin{lstlisting}
// create new document
string path = @"C:\Data\sample.xlsx";
// add title row 
List<string> titleRow = new List<string>();
titleRow.Add("This is the 1st cell");
titleRow.Add("This is the 2nd cell");
openXmlExcel.addTitleRow(titleRow);
// add scatter chart
List<double[]> xValues = new List<double[]>();
List<double[]> yValues = new List<double[]>();
List<string> lineNames = new List<string>();
// first line
xValues.Add(new double[5] { -2, -1, 0, 1, 2 });
yValues.Add(new double[5] { -2, -1, 0, 1, 2 });
lineNames.Add("Line 1");
// second line
xValues.Add(new double[5] { -2, -1, 0, 1, 2 });
yValues.Add(new double[5] { -3, -1, 0, 1, 3 });
lineNames.Add("Line 2");
// save document 
openXmlExcel.saveFile();
\end{lstlisting}

\newpage

\section{Verzeichnisse}
% Bildverzeichnis:
\listoffigures 
%
% Abkürzungsverzeichnis:
%Abkürzungsverzeichnis: Alle Abkürzungen, die bekannt sind finden sich in diesem Verzeichnis.
%Nur diejenigen, die im Dokument auch verwendet werden kommen in die Liste der Compilierung 
%
%der Name der Überschrift kann frei gewählt werden (in diesem Fall: Abkürzungen):
%
%\section*{Acronymverzeichnis}					%Alternative Überschrift für das Verzeichnis
\section*{Abkürzungen}
%
%usage im Dokument: \acs{DA}
%usage im Dokument: \acs{Kurzzeichen}
\begin{acronym}[längstesAcronym]
\setlength{\itemsep}{-0.3 \parsep}
%
% Einträge für Abkürzungen:
%usage im Verzeichnis: \acro{Kurzzeichen}[Inhalt der im Text steht]{Text der als Erklärung im Glossar steht!}
%
\acro{DA}[Diplomarbeit]{Text für die Erklärung der Diplomarbeit}
\acro{ESD}[ESD]{Electrostatic Discharge}
\acro{EMC}[EMC]{Electromagnetic Compatibility}
%
%
% weitere Einträge ...
%
\end{acronym}

%
% Literaturverzeichnis:
%Literaturverzeichnis: Alle Zitate und Veröffentlichungen, die bekannt sind finden sich in diesem Verzeichnis.
%Nur diejenigen, die im Dokument auch verwendet werden kommen in die Liste der Compilierung 
%
%usage im Dokument: \cite{Kurzzeichen}   hier wird im Dokument nur der Text in eckiger Klammer abgebildet
%usage im Dokument: \cite[S~105]{11_Killinger} hier wird auch eine Seitenanzahl dazugestellt
\begin{thebibliography}{22}
%
% Einträge für Literaturverzeichnis:
%usage im Verzeichnis: \bibitem[Werk, Jahrgang oder Auflage]{Kurzzeichen}{Text der als Erklärung im Glossar steht!}
%
   \bibitem[Killinger, 1998]{11_Killinger} Sprache heute : Schriftverkehr; Deutsch für berufsbildende Schulen,\newline 
Seit der Einführung der Bildungsstandards steht der Erwerb fachlicher und sozialer Kompetenzen für Alltag und Beruf im Fokus des Unterrichts.\newline
Autoren: Walter Pirnath, Robert Killinger, Josef Neumüller, 1998
%
\bibitem[vgl. Mustermann, 2020]{12_MM} Skriptum: Test im Zuammenhang mit  ...\newline 
Autoren: Max Mustermann, 2020
%
\bibitem[Werk 1, Jahrgang oder Auflage 1]{BlindzitatKurzzeichen_1}{Text der als Erklärung im Literaturverzeichnis steht!}
%
% weitere Einträge ...
%
\end{thebibliography}
				%ACHTUNG: für die Gesamtcompilierung auskommentieren!!
%
\end{document}